\sezione{8}{Documentare uno schema}{slide 146–150}

\begin{inbreve}
Uno schema E-R non basta da solo: va accompagnato da un \hl{dizionario dei dati}, che descrive a parole entità e relationship, e dall'elenco dei \hl{vincoli non esprimibili} con i costrutti del modello.
\end{inbreve}

\subsection{Lo schema completo dell'università}

\par\noindent\begin{minipage}{\linewidth}
\begin{center}\schemauniversita{true}\end{center}
\didascalia{Schema dell'università con cardinalità e identificatori (slide 147). Dipartimento è identificato esternamente da Nome e Università.}
\end{minipage}\par

\begin{attenzione}[Due versioni dello stesso schema]
Le slide mostrano lo schema due volte con cardinalità diverse. Nella slide 111 Direzione ha $(0,1)$ dal lato di Dipartimento e Afferenza $(1,1)$ dal lato di Docente; nella slide 147 Direzione è $(1,1)$ per Dipartimento (ogni dipartimento ha un direttore) e Afferenza è $(0,1)$ per Docente (possono esistere docenti che non afferiscono a nessun dipartimento, come richiede il vincolo 4 qui sotto). Qui si usa la versione della slide 147.
\end{attenzione}

\subsection{Dizionario dei dati}

\begin{center}
\small
\renewcommand{\arraystretch}{1.35}\renewcommand\tabularxcolumn[1]{m{#1}}
\rowcolors{2}{chiaro}{white}
\begin{tabularx}{\linewidth}{>{\raggedright\headfont\bfseries\small}m{2.5cm} >{\raggedright\arraybackslash}X >{\raggedright\arraybackslash}m{3.6cm} >{\raggedright\arraybackslash}m{2.8cm}}
\rowcolor{navy}\thd{Entità} & \thd{Descrizione} & \thd{Attributi} & \thd{Identificatore}\\
Docente & Docente dell'università & Codice, Cognome, Stipendio & Codice\\
Progetto & Progetti di ricerca & Nome, Budget & Nome\\
Dipartimento & Struttura dell'ateneo & Nome, Telefono & Nome, Università\\
Università & Ateneo & Nome, Città, Indirizzo & Nome\\
\end{tabularx}
\end{center}
\didascalia{Dizionario dei dati: entità (slide 148). Stipendio compare nel dizionario ma non nel disegno: il dizionario può aggiungere dettagli.}

\begin{center}
\small
\renewcommand{\arraystretch}{1.35}\renewcommand\tabularxcolumn[1]{m{#1}}
\rowcolors{2}{chiaro}{white}
\begin{tabularx}{\linewidth}{>{\raggedright\headfont\bfseries\small}m{2.8cm} >{\raggedright\arraybackslash}X >{\raggedright\arraybackslash}m{4.2cm} >{\raggedright\arraybackslash}m{1.6cm}}
\rowcolor{navy}\thd{Relationship} & \thd{Descrizione} & \thd{Componenti} & \thd{Attributi}\\
Direzione & Direzione di un dipartimento & Docente, Dipartimento & \\
Afferenza & Afferenza a un dipartimento & Docente, Dipartimento & Data\\
Partecipazione & Partecipazione a un progetto & Docente, Progetto & \\
Composizione & Composizione dell'ateneo & Dipartimento, Università & \\
\end{tabularx}
\end{center}
\didascalia{Dizionario dei dati: relationship (slide 149).}

\subsection{Vincoli non esprimibili}

\begin{tecnica}[Vincoli di integrità sui dati]
Alcune regole della realtà non si possono disegnare con entità, relationship, cardinalità, identificatori e generalizzazioni. Si elencano a parole:
\begin{passi}
  \item Il direttore di un dipartimento deve afferire a tale dipartimento.
  \item Un docente non deve avere uno stipendio maggiore del direttore del dipartimento al quale afferisce.
  \item Un dipartimento di un'università con sede a Roma deve essere diretto da un docente con più di dieci anni di anzianità.
  \item Un docente che non afferisce a nessun dipartimento non può partecipare a nessun progetto.
\end{passi}
\end{tecnica}

\begin{intuizione}[Perché non si possono disegnare]
Le cardinalità parlano di \hl{una relationship alla volta}. Il vincolo 1 lega \emph{due} relationship (Direzione e Afferenza devono «puntare» allo stesso dipartimento); il vincolo 2 confronta \emph{valori} di attributi; il vincolo 3 dipende da un valore (la città); il vincolo 4 condiziona una relationship alla partecipazione a un'altra. Nessun costrutto del modello esprime questo tipo di regole.
\end{intuizione}

\begin{esercizi}
\begin{passi}
  \item Aggiungi al dizionario dei dati una riga per l'entità \emph{Sede} dello schema A del capitolo 6.
  \item Scrivi un vincolo non esprimibile per lo schema della biblioteca (capitolo 9).
\end{passi}
\end{esercizi}

\begin{appunti}{Riepilogo del capitolo}
\pag{146} Documentazione: dizionario dei dati (entità e relationship) + vincoli non esprimibili.
\pag{148} Dizionario delle entità: descrizione, attributi, identificatore.
\pag{149} Dizionario delle relationship: descrizione, componenti, attributi.
\pag{150} Vincoli non esprimibili: regole fra più costrutti o sui valori.
\end{appunti}
